\documentclass[10pt,a4paper]{amsart}
\usepackage{amsmath,amssymb,amsthm,graphicx}
\usepackage{verbatim}
\newenvironment{reference}{\comment}{\endcomment}
%\newenvironment{reference}{}{}
\newenvironment{slogan}{\comment}{\endcomment}
\newenvironment{history}{\comment}{\endcomment}
% For commutative diagrams we use Xy-pic
\usepackage[all]{xy}

% We use 2cell for 2-commutative diagrams.
\xyoption{2cell}
% Theorem environments.
%
\theoremstyle{plain}
\newtheorem{theorem}[subsection]{Theorem}
\newtheorem{proposition}[subsection]{Proposition}
\newtheorem{lemma}[subsection]{Lemma}

\theoremstyle{definition}
\newtheorem{definition}[subsection]{Definition}
\newtheorem{example}[subsection]{Example}
\newtheorem{exercise}[subsection]{Exercise}
\newtheorem{situation}[subsection]{Situation}

\theoremstyle{remark}
\newtheorem{remark}[subsection]{Remark}
\newtheorem{remarks}[subsection]{Remarks}

\numberwithin{equation}{subsection}

% Macros
%
\def\lim{\mathop{\mathrm{lim}}\nolimits}
\def\colim{\mathop{\mathrm{colim}}\nolimits}
\def\Spec{\mathop{\mathrm{Spec}}}
\def\Hom{\mathop{\mathrm{Hom}}\nolimits}
\def\Ext{\mathop{\mathrm{Ext}}\nolimits}
\def\SheafHom{\mathop{\mathcal{H}\!\mathit{om}}\nolimits}
\def\SheafExt{\mathop{\mathcal{E}\!\mathit{xt}}\nolimits}
\def\Sch{\mathit{Sch}}
\def\Mor{\mathop{\mathrm{Mor}}\nolimits}
\def\Ob{\mathop{\mathrm{Ob}}\nolimits}
\def\Sh{\mathop{\mathit{Sh}}\nolimits}
\def\NL{\mathop{N\!L}\nolimits}
\def\CH{\mathop{\mathrm{CH}}\nolimits}
\def\proetale{{pro\text{-}\acute{e}tale}}
\def\etale{{\acute{e}tale}}
\def\QCoh{\mathit{QCoh}}
\def\Ker{\mathop{\mathrm{Ker}}}
\def\Im{\mathop{\mathrm{Im}}}
\def\Coker{\mathop{\mathrm{Coker}}}
\def\Coim{\mathop{\mathrm{Coim}}}

% Boxtimes
%
\DeclareMathSymbol{\boxtimes}{\mathbin}{AMSa}{"02}

%
% Macros for moduli stacks/spaces
%
\def\QCohstack{\mathcal{QC}\!\mathit{oh}}
\def\Cohstack{\mathcal{C}\!\mathit{oh}}
\def\Spacesstack{\mathcal{S}\!\mathit{paces}}
\def\Quotfunctor{\mathrm{Quot}}
\def\Hilbfunctor{\mathrm{Hilb}}
\def\Curvesstack{\mathcal{C}\!\mathit{urves}}
\def\Polarizedstack{\mathcal{P}\!\mathit{olarized}}
\def\Complexesstack{\mathcal{C}\!\mathit{omplexes}}
% \Pic is the operator that assigns to X its picard group, usage \Pic(X)
% \Picardstack_{X/B} denotes the Picard stack of X over B
% \Picardfunctor_{X/B} denotes the Picard functor of X over B
\def\Pic{\mathop{\mathrm{Pic}}\nolimits}
\def\Picardstack{\mathcal{P}\!\mathit{ic}}
\def\Picardfunctor{\mathrm{Pic}}
\def\Deformationcategory{\mathcal{D}\!\mathit{ef}}
\usepackage{cleveref}
\crefname{originallemma}{original lemma}{original lemmas}
\Crefname{originallemma}{Original lemma}{Original lemmas}
\makeatletter
\newlabel{lemma-Noetherian}{{lemma-Noetherian}{source}{}{sourcepointer-1}{}}
\newlabel{lemma-Noetherian@cref}{{[originallemma][0][]lemma-Noetherian}{[1][0][]source}{}{}{}}
\makeatother
\crefname{originallemma}{original lemma}{original lemmas}
\Crefname{originallemma}{Original lemma}{Original lemmas}
\makeatletter
\newlabel{lemma-constructible-Noetherian-space}{{lemma-constructible-Noetherian-space}{source}{}{sourcepointer-2}{}}
\newlabel{lemma-constructible-Noetherian-space@cref}{{[originallemma][0][]lemma-constructible-Noetherian-space}{[1][0][]source}{}{}{}}
\makeatother
\crefname{originallemma}{original lemma}{original lemmas}
\Crefname{originallemma}{Original lemma}{Original lemmas}
\makeatletter
\newlabel{lemma-constructible-map-Noetherian}{{lemma-constructible-map-Noetherian}{source}{}{sourcepointer-3}{}}
\newlabel{lemma-constructible-map-Noetherian@cref}{{[originallemma][0][]lemma-constructible-map-Noetherian}{[1][0][]source}{}{}{}}
\makeatother
\title{Characterization of constructible subsets of a Noetherian space by density conditions}
\author{The Stacks Project Authors}
\date{stacks-project:topology.tex:local-snapshot-20261012; GNU FDL 1.2 or later}
\begin{document}
\maketitle
\noindent\textit{Editorial note.} All delivered mathematical text is the author's own, copied from the pinned Stacks Project chapter source sources/stacks-topology/pinned-original-source.tex; no statement or proof step was rewritten, reordered or removed.\par\medskip
\noindent\textit{Editorial note.} Preamble provenance: the Stacks Project keeps one single shared preamble for every chapter; this card takes it directly from the staged sources/stacks-topology/preamble.tex as three verbatim spans (lines 8-12, 14-18 and 154-217), so no derived file was written and nothing was added or redefined. The wrapper class and the base packages are supplied by the assembler.\par\medskip
\noindent\textit{Editorial note.} Bibliography provenance: neither the counted statement nor its proof contains a citation (this lane rescans the whole span before writing the card), so the card carries no bibliography and no entry was reconstructed.\par\medskip
\noindent\textit{Editorial note.} Selection provenance: this card was opened under the file-based control contract (CONTROL.md), with the mechanical pre-card gate run on its statement and proof interval before the card was written.\par\medskip
\noindent\textit{Allowed standard prerequisites (editorial).}\begin{itemize}
\item The statement that a constructible subset of a Noetherian space is a finite union of locally closed subsets. Invoked inside the counted proof (source lines 3164-3170 and 3187-3192). Turns constructibility into a finite union of locally closed pieces.
\item The decomposition of a Noetherian space into finitely many irreducible components. Invoked inside the counted proof (source lines 3176-3178). Splits the minimal counterexample into components.
\end{itemize}\medskip
\ifdefined\hypertarget\hypertarget{sourcepointer-1}{}\fi
\noindent\textit{Source reference.} Original lemma lemma-Noetherian: sources/stacks-topology/pinned-original-source.tex, source lines 1470--1470. Navigation pointer only: the counted text refers to the Stacks item lemma-Noetherian, which is an existing labeled result of this same chapter and is not part of this card. The pointer carries the target's own label as its number and its definition line as its span; its content is not reproduced here.\par\medskip
\ifdefined\hypertarget\hypertarget{sourcepointer-2}{}\fi
\noindent\textit{Source reference.} Original lemma lemma-constructible-Noetherian-space: sources/stacks-topology/pinned-original-source.tex, source lines 3088--3088. Navigation pointer only: the counted text refers to the Stacks item lemma-constructible-Noetherian-space, which is an existing labeled result of this same chapter and is not part of this card. The pointer carries the target's own label as its number and its definition line as its span; its content is not reproduced here.\par\medskip
\ifdefined\hypertarget\hypertarget{sourcepointer-3}{}\fi
\noindent\textit{Source reference.} Original lemma lemma-constructible-map-Noetherian: sources/stacks-topology/pinned-original-source.tex, source lines 3100--3100. Navigation pointer only: the counted text refers to the Stacks item lemma-constructible-map-Noetherian, which is an existing labeled result of this same chapter and is not part of this card. The pointer carries the target's own label as its number and its definition line as its span; its content is not reproduced here.\par\medskip
\begin{lemma}
\label{lemma-characterize-constructible-Noetherian}
Let $X$ be a Noetherian topological space.
Let $E \subset X$ be a subset.
The following are equivalent:
\begin{enumerate}
\item $E$ is constructible in $X$, and
\item for every irreducible closed $Z \subset X$ the intersection
$E \cap Z$ either contains a nonempty open of $Z$ or is not dense in $Z$.
\end{enumerate}
\end{lemma}

\begin{proof}
Assume $E$ is constructible and $Z \subset X$ irreducible closed.
Then $E \cap Z$ is constructible in $Z$ by
Lemma \ref{lemma-constructible-map-Noetherian}.
Hence $E \cap Z$ is a finite union of nonempty locally closed subsets
$T_i$ of $Z$. Clearly if none of the $T_i$ is open in $Z$, then
$E \cap Z$ is not dense in $Z$. In this way we see that (1) implies (2).

\medskip\noindent
Conversely, assume (2) holds. Consider the set $\mathcal{S}$ of closed
subsets $Y$ of $X$ such that $E \cap Y$ is not constructible in $Y$.
If $\mathcal{S} \not = \emptyset$, then it has a smallest element $Y$
as $X$ is Noetherian.
Let $Y = Y_1 \cup \ldots \cup Y_r$ be the decomposition of $Y$ into its
irreducible components, see
Lemma \ref{lemma-Noetherian}.
If $r > 1$, then each $Y_i \cap E$ is constructible in $Y_i$ and hence
a finite union of locally closed subsets of $Y_i$. Thus $E \cap Y$
is a finite union of locally closed subsets of $Y$ too and we conclude
that $E \cap Y$ is constructible in $Y$ by
Lemma \ref{lemma-constructible-Noetherian-space}.
This is a contradiction and so $r = 1$. If $r = 1$, then $Y$ is
irreducible, and by assumption (2) we see that $E \cap Y$ either
(a) contains an open $V$ of $Y$ or (b) is not dense in $Y$.
In case (a) we see, by minimality of $Y$, that $E \cap (Y \setminus V)$
is a finite union of locally closed subsets of $Y \setminus V$. Thus
$E \cap Y$ is a finite union of locally closed subsets of $Y$ and is
constructible by
Lemma \ref{lemma-constructible-Noetherian-space}.
This is a contradiction and so we must be in case (b).
In case (b) we see that $E \cap Y = E \cap Y'$ for some proper closed
subset $Y' \subset Y$. By minimality of $Y$ we see that
$E \cap Y'$ is a finite union of locally closed subsets of $Y'$ and
we see that $E \cap Y' = E \cap Y$ is a finite union of locally closed
subsets of $Y$ and is constructible by
Lemma \ref{lemma-constructible-Noetherian-space}.
This contradiction finishes the proof of the lemma.
\end{proof}

\end{document}
